\section{The finite-dimensional correlation ceiling}\label{sec:qa}
Let $\mathcal A_X,\mathcal A_Y$ be disjoint quasi-local spin
algebras, each the norm closure of finite-site matrix algebras, with
finite local site dimension. A state may be mixed or an arbitrary
infinite-volume state. Measurement families are finite and their
effects belong to these quasi-local algebras.

\begin{proposition}[Quasi-local correlations are approximable]\label{prop:qa}
Every such correlation belongs to $C_{qa}=\overline{C_q}$, the
closure of finite-dimensional tensor-product correlations. If
$\omega_{qa}(G)<1$, its loss in $G$ is at least
$1-\omega_{qa}(G)$. In a legal referee cell this implies
\[
 e\geq\frac{1-\omega_{qa}(G)}{C_kM}.
\]
\end{proposition}
\begin{proof}
Choose increasing finite subsets of $X$ and conditional expectations
onto their matrix algebras, using normalized trace on the omitted
sites. These maps are unital completely positive, preserve a POVM's
positivity and sum, and converge in norm on every quasi-local
observable. Do the same independently in $Y$. Restrict the state to
the resulting finite tensor product of matrix algebras; it is a
density matrix. Thus every approximating correlation is in $C_q$.
Norm convergence of the finitely many effects gives convergence of
all joint probabilities and proves membership in its closure. A Bell
functional is continuous on this finite probability space, proving
the loss floor. Compose the floor with extraction for the energy
inequality.
\end{proof}

This standard approximation principle \cite{scholz} does not depend
on a pure-state infinite-volume tensor splitting. General commuting
operator correlations need not admit these quasi-local measurements;
the distinction between $C_{qa}$ and $C_{qc}$ is substantive
\cite{mip}. The cell energy floor retains its normalization $C_kM$;
it need not be bounded away from zero as the clock grows.
